%% Appendix B: API Reference
\chapter{API Reference}
\label{app:api}

All API endpoints use JSON request/response bodies. Base URL: \texttt{http://SERVER\_IP:5000}

\section{Authentication Endpoints}

\subsection{POST /api/v1/auth/signup}
Creates a new user account.

\textbf{Request Body:}
\begin{tcolorbox}[codebox, title={Request}]
\begin{lstlisting}[language=bash]
{
  "enrollmentNumber": "ENR2026011",  // Unique student ID
  "fullName": "Riya Kapoor",
  "password": "Student@2026!",       // Min 8 chars
  "pin": "123456",                   // 6-digit numeric
  "role": "STUDENT"                  // STUDENT | ADMIN | SUPERVISOR
}
\end{lstlisting}
\end{tcolorbox}

\textbf{Response (201):}
\begin{tcolorbox}[codebox, title={Response}]
\begin{lstlisting}[language=bash]
{
  "token": "eyJhbGciOiJIUzI1NiJ9...",
  "user": {
    "id": "uuid",
    "enrollmentNumber": "ENR2026011",
    "fullName": "Riya Kapoor",
    "role": "STUDENT"
  }
}
\end{lstlisting}
\end{tcolorbox}

\subsection{POST /api/v1/auth/login}
Authenticates a user and returns a JWT token.

\textbf{Request Body:}
\begin{tcolorbox}[codebox, title={Request}]
\begin{lstlisting}[language=bash]
{
  "enrollmentNumber": "ADMIN01",
  "password": "Admin@ALAMS2026!"
}
\end{lstlisting}
\end{tcolorbox}

\section{Admin Endpoints (JWT Required)}

All admin endpoints require the header:
\begin{tcolorbox}[codebox, title={Authorization Header}]
\begin{lstlisting}[language=bash]
Authorization: Bearer <jwt_token>
\end{lstlisting}
\end{tcolorbox}

\subsection{GET /api/v1/admin/analytics/pilot}
Returns computed pilot KPI metrics.

\textbf{Response (200):}
\begin{tcolorbox}[codebox, title={Response}]
\begin{lstlisting}[language=bash]
{
  "loginSuccessRate": 97.5,
  "successfulLogins": 39,
  "failedLoginAlerts": 1,
  "totalAttempts": 40,
  "avgQrLatencyMs": 3821,
  "avgPinLatencyMs": 1450,
  "qrSampleCount": 30,
  "pinSampleCount": 9,
  "hourlyDistribution": [
    {"hour": 9, "count": 5},
    {"hour": 10, "count": 12},
    ...
  ],
  "verificationBreakdown": {
    "QR_CODE": 30,
    "PIN_FALLBACK": 9,
    "ADMIN_OVERRIDE": 0,
    "total": 39
  },
  "watchdogEnforcementCount": 0,
  "watchdogActiveCount": 8,
  "clientOnlineCount": 10,
  "totalComputers": 10,
  "pilotTarget": 10,
  "deploymentProgress": 10
}
\end{lstlisting}
\end{tcolorbox}

\subsection{POST /api/v1/admin/computers/approve}
Approves a pending workstation registration.

\textbf{Request Body:}
\begin{tcolorbox}[codebox, title={Request}]
\begin{lstlisting}[language=bash]
{
  "computerId": "uuid-of-pending-computer",
  "pcNumber": "PC-A01",
  "labId": "uuid-of-lab",
  "fallbackEnabled": true
}
\end{lstlisting}
\end{tcolorbox}

\subsection{POST /api/v1/admin/computers/remote-unlock}
\textbf{Request:} \verb|{ "computerId": "uuid" }|

\subsection{POST /api/v1/admin/computers/remote-lock}
\textbf{Request:} \verb|{ "computerId": "uuid" }|

\section{Client (Workstation) Endpoints}

\subsection{GET /api/v1/client/qr-token}
Generates a signed QR token for the given computer.

\textbf{Query Parameter:} \verb|?computerId=uuid|

\textbf{Response:}
\begin{tcolorbox}[codebox, title={Response}]
\begin{lstlisting}[language=bash]
{
  "token": "eyJhbGciOiJIUzI1NiJ9..."  // 60-second JWT
}
\end{lstlisting}
\end{tcolorbox}

\subsection{POST /api/v1/client/fallback-auth}
Authenticates a student with enrollment number and PIN.

\textbf{Request:}
\begin{tcolorbox}[codebox, title={Request}]
\begin{lstlisting}[language=bash]
{
  "enrollmentNumber": "ENR2026001",
  "pin": "123456",
  "computerId": "uuid"
}
\end{lstlisting}
\end{tcolorbox}

\subsection{POST /api/v1/client/watchdog-heartbeat}
Reports watchdog liveness. Called every 10 seconds.

\textbf{Request:} \verb|{ "computerId": "uuid" }|

\subsection{POST /api/v1/client/watchdog-alert}
Reports a security bypass event.

\textbf{Request:}
\begin{tcolorbox}[codebox, title={Request}]
\begin{lstlisting}[language=bash]
{
  "computerId": "uuid",
  "alertType": "watchdog_kill",
  "severity": "CRITICAL",
  "details": "AlamsClient.exe was terminated while desktop was active."
}
\end{lstlisting}
\end{tcolorbox}

\section{Mobile Gateway Endpoint}

\subsection{POST /api/v1/mobile/verify-unlock}
Verifies a QR token and unlocks the workstation. Requires JWT.

\textbf{Request:}
\begin{tcolorbox}[codebox, title={Request}]
\begin{lstlisting}[language=bash]
{
  "qrToken": "eyJhbGciOiJIUzI1NiJ9..."
}
\end{lstlisting}
\end{tcolorbox}

\textbf{Response (200):}
\begin{tcolorbox}[codebox, title={Response}]
\begin{lstlisting}[language=bash]
{
  "message": "Workstation unlocked successfully",
  "computer": {
    "pcNumber": "PC-A01",
    "deviceName": "ALAMS-TURING-01"
  },
  "sessionId": "uuid"
}
\end{lstlisting}
\end{tcolorbox}

\section{WebSocket Protocol}

Connection URL: \texttt{ws://SERVER\_IP:5000}

\subsection{Client → Server Messages}

\begin{table}[h!]
\centering
\caption{WebSocket Messages: Client to Server}
\begin{tabularx}{\textwidth}{lX}
\toprule
\textbf{Message Type} & \textbf{Payload}\\
\midrule
\texttt{register} & \verb|{type, fingerprint, deviceName, macAddress, ipAddress}|\\
\rowcolor{aurxonlightgray}
\texttt{heartbeat} & \verb|{type, status: "locked"|"in_use"}|\\
\bottomrule
\end{tabularx}
\end{table}

\subsection{Server → Client Messages}

\begin{table}[h!]
\centering
\caption{WebSocket Messages: Server to Client}
\begin{tabularx}{\textwidth}{lX}
\toprule
\textbf{Message Type} & \textbf{Payload}\\
\midrule
\texttt{approved} & \verb|{type, computerId, pcNumber, qrSeed, fallbackEnabled}|\\
\rowcolor{aurxonlightgray}
\texttt{unlock} & \verb|{type, enrollmentNumber, sessionId}|\\
\texttt{lock} & \verb|{type}|\\
\bottomrule
\end{tabularx}
\end{table}
